\documentclass[10pt]{article}

\usepackage[T1]{fontenc}
\usepackage[utf8]{inputenc}
\usepackage[margin=0.78in]{geometry}
\usepackage{amsmath,amssymb,amsthm,mathtools}
\usepackage{booktabs}
\usepackage{microtype}
\usepackage{xcolor}
\usepackage{hyperref}
\definecolor{linkblue}{HTML}{1D4E89}

\hypersetup{
  pdftitle={Deterministic Certification for the Bounded Normal Mean: An Adaptive-Support Computation with an Independent Verifier},
  pdfauthor={Maarten Linus Herrmann (ORCID: 0009-0005-8721-6588)},
  pdfsubject={Working manuscript, 27 October 2026},
  pdfkeywords={bounded normal mean, minimax estimation, least favorable prior, interval arithmetic, branch-and-bound, computer-assisted proof},
  pdflang={en},
  colorlinks=true,
  linkcolor=black,
  citecolor=linkblue,
  urlcolor=linkblue
}

\setlength{\parskip}{2.5pt}
\setlength{\parindent}{1.2em}
\setlength{\abovedisplayskip}{5pt plus 2pt minus 2pt}
\setlength{\belowdisplayskip}{5pt plus 2pt minus 2pt}
\raggedbottom

\newtheorem{assumption}{Assumption}
\newtheorem{definition}{Definition}
\newtheorem{lemma}{Lemma}
\newtheorem{proposition}{Proposition}
\newtheorem{theorem}{Theorem}
\newtheorem{corollary}{Corollary}
\newtheorem{remark}{Remark}

\newcommand{\R}{\mathbb R}
\newcommand{\E}{\mathbb E}
\newcommand{\Var}{\operatorname{Var}}
\newcommand{\supp}{\operatorname{supp}}
\newcommand{\Normal}{\mathsf N}
\newcommand{\phiN}{\phi}
\newcommand{\risk}{r}
\newcommand{\BR}{b}
\newcommand{\vstar}{v^\star}
\newcommand{\eps}{\varepsilon}
\newcommand{\code}[1]{\texttt{#1}}

\title{\textbf{Deterministic Certification for the Bounded Normal Mean}\\
\large An Adaptive-Support Computation with an Independent Verifier}
\author{
Maarten Linus Herrmann\\
\small Erasmus School of Economics\\
\small Erasmus University Rotterdam\\
\small ORCID: \href{https://orcid.org/0009-0005-8721-6588}{0009-0005-8721-6588}
}
\date{Working manuscript \textbullet{} 27 October 2026}

\usepackage{xurl}
\begin{document}
\maketitle

\begin{abstract}
We consider estimation of a normal mean in a bounded interval under squared
error. A finite prior gives a lower bound on minimax risk through its Bayes
risk and an upper bound through the maximum risk of its Bayes rule. We use
Kempthorne's established exchange strategy to find a candidate symmetric prior
and an independent outward-rounded program to verify the resulting bracket.
The verifier reconstructs Gaussian-mixture ratios, bounds Simpson quadrature
remainders and Gaussian tails, and covers the continuous parameter interval
by branch-and-bound. It does not rely on the discovery solver's optimizer or
branch tree. For the interval $[-4,4]$, it certifies minimax risk in
$[0.8150887565050,0.8150887615302]$, with width below
$5.026\times10^{-9}$. Posterior-flow identities provide a curvature
bound for the discovery calculation; an ablation shows that its specialized
common-mean enclosure yields a modest reduction in branch nodes for this
instance. The certificate is a computer-assisted result conditional on the
correctness of the verifier and its arithmetic implementation. We do not
claim a new exchange algorithm or a finite-time convergence guarantee for
the discovery procedure.
\end{abstract}

\noindent\textbf{Keywords:} bounded normal mean; minimax estimation; least
favorable prior; interval arithmetic; branch-and-bound; computer-assisted
proof.

\noindent\textbf{JEL codes:} C14, C44, C63.
\qquad
\textbf{MSC 2020:} 62C20, 65G20, 90C26.

\section{Introduction}
\label{sec:introduction}

A worst-case risk guarantee for the bounded-normal mean requires a global
calculation over a continuum. Let
\[
  Y=\theta+Z,\qquad Z\sim\Normal(0,1),\qquad \theta\in[-m,m],
\]
and use loss $(d(Y)-\theta)^2$.  The objective is the minimax value
\[
  \vstar=\inf_d\sup_{\theta\in[-m,m]}
    \E_\theta\{d(Y)-\theta\}^2 .
\]

Montiel Olea and Zubova \cite{montielzubova2026} recently gave a stochastic
mirror-ascent method with an ex ante equally spaced parameter grid and
finite-sample probability guarantees.  Their paper prompted the present
project.  Our question is different: after a candidate finitely supported
prior has been found, can a deterministic program certify its Bayes risk and
the global risk of its Bayes rule without a parameter grid or Monte Carlo?

For the $m=4$ instance, the final
artifact contains three positive support levels, outward-rounded
quadrature and tail enclosures, and a continuous branch-and-bound proof.
A second program reconstructs the calculation from the compact certificate
without importing the discovery solver.

\paragraph{Roadmap.}
Section~\ref{sec:lineage} separates inherited theory, independent derivations,
and implementation.  Section~\ref{sec:formulation} proves the exact
directional derivative and minimax bracket.  Section~\ref{sec:flow} develops
the posterior-flow representation.  Sections~\ref{sec:separation} and
\ref{sec:algorithm} give the continuous separator and adaptive procedure.
Section~\ref{sec:verifier} describes the independent verifier.
Sections~\ref{sec:results}--\ref{sec:comparison} report tests and compare with
Montiel Olea and Zubova.  The appendices supply all proof algebra, numerical
error bounds, lineage details, and exact reproduction commands.

\paragraph{Contribution and lineage.}
Kempthorne \cite{kempthorne1987} supplies adaptive atom insertion, joint
optimization of weights and locations, pruning, the directional derivative,
and the Bayes-risk/maximum-risk stopping bracket.  Interval arithmetic,
dependency-aware enclosures and semiconvex branch-and-bound provide the
numerical foundations.  This paper contributes a specialized deterministic
bounded-normal implementation, a continuous separator, and an independent
outward-rounded certificate.  Lemma~\ref{lem:common-mean-bound} gives a
posterior-moment enclosure derived for this implementation; the ablation shows
a four-node reduction at $m=4$, so it is treated as a supporting device rather
than the principal result.

\section{Exact contribution and relation to prior work}
\label{sec:lineage}

This section states lineage before technical results become operative.
Table~\ref{tab:lineage} records the status used throughout the paper.

\begin{table}[ht]
\centering
\caption{Technical lineage and the role of each ingredient.}
\label{tab:lineage}
\small
\begin{tabular}{p{0.28\linewidth}p{0.22\linewidth}p{0.42\linewidth}}
\toprule
Object & Role & Source and exact relation \\
\midrule
Adaptive support insertion; reoptimization of masses and locations; zero-mass
or coincident-point reduction & Prior art &
Kempthorne \cite[pp.~175--177]{kempthorne1987}, Steps 1--4 and discussion.\\
Directional derivative and equalizer/stopping condition & Prior art &
Kempthorne \cite[Theorems 2.1--2.2, pp.~172--174]{kempthorne1987}; the exact
Gaussian algebra is supplied here for auditability.\\
Finite-support structure & Prior art &
Kempthorne \cite[Theorem 2.4, p.~175]{kempthorne1987}, with later structural
work including \cite{dytso2018,dytso2022}.\\
Deterministic global optimization of fixed finite prior parameterizations &
Prior art &
Gourdin, Jaumard, and MacGibbon \cite[Sections 2--3]{gourdin1994}.\\
Stochastic grid-based approximation and Wasserstein/estimator convergence &
Related recent method &
Theorems 1--4 of \cite{montielzubova2026}.\\
Posterior-flow identities $\mu'=V$, $V'=\kappa_3$ & Standard exponential-family
algebra & Re-derived in Proposition~\ref{prop:flow}.\\
Sharp support bound for $\kappa_3$ & Prior art &
Equivalent to Sharma, Gupta, and Kapoor
\cite[Theorem 2.2 and (2.20)--(2.21)]{sharma2010}.\\
Common-mean product enclosure & Supporting derivation &
Lemma~\ref{lem:common-mean-bound}; a posterior-moment enclosure with a
four-node reduction at $m=4$.\\
Arb verifier and $m=4$ certificate & Release contribution &
Separate direct-ratio formula, Simpson remainder, tail bound, and
branch-and-bound; source and certificate accompany the paper.\\
\bottomrule
\end{tabular}
\end{table}

Casella and Strawderman \cite{casella1981} establish foundational exact
results for small bounds.  Feldman and Brown \cite{feldman1989} and Donoho,
Liu, and MacGibbon \cite{donoho1990} develop numerical and risk-bound
calculations.  Gourdin et al.\ \cite{gourdin1994} use theoretically convergent
Lipschitz and decomposition global optimization for finite-dimensional
parameterizations.  Kent et al.\ \cite{kent2021} give a modified
Frank--Wolfe method in probability space with oracle assumptions, while
Guggenberger and Huang \cite{guggenberger2025} use fictitious play and
grid-based best responses in a different minimax-regret problem.  Aradillas
Fernández et al.\ \cite{aradillas2026} prove epsilon-minimax guarantees for
mixtures over finitely many decision rules under a worst-case-risk oracle.
The present release connects these lines of work to an outward-rounded
continuous separator and an independently executable certificate for the
scalar bounded-normal problem.

Interval arithmetic is classical \cite{moore2009}.  Everitt
\cite{everitt2018} explains dependency over-enclosure and affine arithmetic
in statistical computation, and Maruyama and Matsuda \cite{maruyama2025}
use interval arithmetic in a computer-assisted minimaxity proof for a fixed
half-Cauchy Bayes estimator.  These works provide the broader
dependency-aware and computer-assisted statistical context.

\section{The bounded-normal minimax problem}
\label{sec:formulation}

This section defines every statistical object and proves the certificate from
the original mixture formulas.

\begin{definition}[Symmetric finite prior]
Let $\pi=\sum_{j=1}^J w_j\delta_{\theta_j}$, where
$\theta_j\in[-m,m]$, $w_j>0$, and $\sum_jw_j=1$.  Symmetry means that
$\pi(A)=\pi(-A)$ for every Borel set $A$.  Define, for $r=0,1,2$,
\[
 f_r(y)=\sum_{j=1}^J w_j\theta_j^r\phiN(y-\theta_j),
 \qquad
 \phiN(x)=(2\pi)^{-1/2}e^{-x^2/2}.
\]
\end{definition}

Since $f_0(y)>0$, the posterior mean and Bayes estimator are
\[
  \delta_\pi(y)=\frac{f_1(y)}{f_0(y)}.
\]
The frequentist risk at $t\in[-m,m]$ and the Bayes risk are
\begin{align}
  \risk(t,\delta_\pi)
   &=\int_{\R}\{\delta_\pi(y)-t\}^2\phiN(y-t)\,dy, \label{eq:risk}\\
  \BR(\pi)
   &=\int_{\R}\left\{f_2(y)-\frac{f_1(y)^2}{f_0(y)}\right\}dy. \label{eq:bayes}
\end{align}
Fubini's theorem gives the identity
\begin{equation}
  \BR(\pi)=\sum_{j=1}^Jw_j\risk(\theta_j,\delta_\pi).
  \label{eq:bayes-average}
\end{equation}
The integrands are dominated because $|\delta_\pi(y)|\le m$.

\begin{proposition}[Exact feasible directional derivative]
\label{prop:direction}
For $t\in[-m,m]$ and $\pi_\alpha=(1-\alpha)\pi+\alpha\delta_t$,
\[
 \left.\frac{d}{d\alpha}\BR(\pi_\alpha)\right|_{\alpha=0+}
 =\risk(t,\delta_\pi)-\BR(\pi).
\]
\end{proposition}

\begin{proof}
Write $h_r(y)=t^r\phiN(y-t)-f_r(y)$, so that
$f_{r,\alpha}=f_r+\alpha h_r$.  Pointwise differentiation of
$f_{2,\alpha}-f_{1,\alpha}^2/f_{0,\alpha}$ gives
\[
 h_2-\frac{2f_1h_1}{f_0}+\frac{f_1^2h_0}{f_0^2}.
\]
Substitute $h_r$ and $\delta_\pi=f_1/f_0$.  The terms involving the original
mixture equal
\[
 -f_2+2\delta_\pi f_1-\delta_\pi^2f_0
 =-f_2+\frac{f_1^2}{f_0},
\]
whereas the atom terms equal
\[
 \{t^2-2t\delta_\pi(y)+\delta_\pi(y)^2\}\phiN(y-t).
\]
Integration yields $\risk(t,\delta_\pi)-\BR(\pi)$.  Differentiation under the
integral is justified in Appendix~\ref{app:proofs}; the key facts are bounded
support and strict positivity of the Gaussian mixture.
\end{proof}

\begin{corollary}[Exact minimax certificate]
\label{cor:certificate}
For every finite prior $\pi$ supported on $[-m,m]$,
\[
 \BR(\pi)\le \vstar\le
 U(\pi):=\sup_{t\in[-m,m]}\risk(t,\delta_\pi),
\]
and therefore
\[
 0\le \vstar-\BR(\pi)
 \le U(\pi)-\BR(\pi)
 =\sup_{t\in[-m,m]}
   \{\risk(t,\delta_\pi)-\BR(\pi)\}.
\]
\end{corollary}

\begin{proof}
For any decision rule $d$, its maximum risk is at least its $\pi$-average
risk, which is at least the Bayes risk $\BR(\pi)$.  Taking the infimum over
$d$ proves the lower bound.  The Bayes rule $\delta_\pi$ is an admissible
competitor in the minimax infimum, which proves the upper bound.  Subtraction
and Proposition~\ref{prop:direction} give the last display.
\end{proof}

Corollary~\ref{cor:certificate}, essentially Kempthorne's stopping condition,
is the logical core of the computation: a lower enclosure for $\BR(\pi)$ and
an upper enclosure for $U(\pi)$ produce a rigorous interval for $\vstar$.

\section{Posterior-flow representation}
\label{sec:flow}

The separator must control a ratio of Gaussian mixtures without replacing
its shared numerator and denominator by independent intervals.  This section
identifies a small posterior state that retains the relevant dependence.

For an observation $y$, let
\[
 p_j(y)=\frac{w_j\exp(\theta_jy-\theta_j^2/2)}
 {\sum_k w_k\exp(\theta_ky-\theta_k^2/2)}.
\]
The omitted common factor $\phiN(y)$ cancels.  Define the posterior mean,
variance, and cumulants
\[
 \mu(y)=\sum_jp_j(y)\theta_j,\quad
 V(y)=\sum_jp_j(y)\{\theta_j-\mu(y)\}^2,\quad
 \kappa_s(y)=\text{$s$th posterior cumulant}.
\]

\begin{proposition}[Posterior flow and risk curvature]
\label{prop:flow}
For a finite bounded prior, all functions above are analytic and
\[
 \mu'=V,\qquad V'=\kappa_3,\qquad \kappa_s'=\kappa_{s+1}.
\]
With $q(t,z)=\mu(t+z)-t$ and $Z\sim\Normal(0,1)$,
\begin{align}
 \risk(t,\delta_\pi)&=\E q(t,Z)^2,\\
 \risk'(t,\delta_\pi)&=2\E\{q(t,Z)(V(t+Z)-1)\},\\
 \risk''(t,\delta_\pi)&=
 2\E\{(V(t+Z)-1)^2+q(t,Z)\kappa_3(t+Z)\}.
 \label{eq:curvature}
\end{align}
\end{proposition}

\begin{proof}
The probabilities form a finite natural exponential family with natural
parameter $y$.  Successive derivatives of
\[
 \log\!\left(\sum_jw_j e^{\theta_jy-\theta_j^2/2}\right)
\]
produce the cumulants.  Further,
$\partial_tq=V(t+z)-1$ and
$\partial_tV=\kappa_3$.  Two differentiations of $\E q^2$ give the displayed
formulas.  Bounded support supplies an integrable uniform bound for each
derivative.
\end{proof}

The only sign-indefinite product in \eqref{eq:curvature} is $q\kappa_3$.
The following representation couples its factors.

\begin{lemma}[Common-mean representation and enclosure]
\label{lem:common-mean-bound}
Let the posterior be supported on $[\alpha,\beta]$.  If $V>0$, define
\[
 \nu=\frac{\E\{\Theta(\Theta-\mu)^2\mid y\}}{V}.
\]
Then $\nu\in[\alpha,\beta]$ and
\[
 \kappa_3=V(\nu-\mu).
\]
If $t\in[a,b]$, $\mu\in M=[\mu_L,\mu_U]$, and $0\le V\le V_U$, then
\begin{align}
 q\kappa_3
 &\le V_U\max_{\mu\in M}
 \bigl\{(\mu-a)_+(\beta-\mu),\,
        (b-\mu)_+(\mu-\alpha)\bigr\}, \label{eq:joint-upper}\\
 q\kappa_3
 &\ge -V_U\max_{\mu\in M}
 \bigl\{(\mu-a)_+(\mu-\alpha),\,
        (b-\mu)_+(\beta-\mu)\bigr\}. \label{eq:joint-lower}
\end{align}
Each maximum is either an endpoint value or the value at the vertex of one
quadratic.
\end{lemma}

\begin{proof}
The weights proportional to
$p_j(\theta_j-\mu)^2$ are nonnegative and sum to $V$, hence $\nu$ is their
mean and lies in the convex hull $[\alpha,\beta]$.  Expanding
$\E\{\Theta(\Theta-\mu)^2\}= \kappa_3+\mu V$ proves the identity.
For a positive product, if $\mu\ge t$ choose the largest feasible
$\nu=\beta$; if $\mu\le t$ choose $\nu=\alpha$.  Maximizing over
$t\in[a,b]$ gives \eqref{eq:joint-upper}.  The opposite choices of $\nu$
give the most negative product and \eqref{eq:joint-lower}.  The case $V=0$
has $\kappa_3=0$ and follows by continuity.
\end{proof}

\paragraph{Equality and degeneracy.}
Equality in the identity holds by construction.  An envelope endpoint can be
attained only when the posterior squared-deviation weights put their mean
$\nu$ at the relevant support endpoint and when $V=V_U$; $V=0$ collapses all
terms to zero.  These conditions are useful diagnostics but are not assumed
by the solver.

\paragraph{Scope of the enclosure.}
The relaxation separates $V$ from $\mu$ through $V_U$ and encloses, rather than
identifies, their exact joint range.  The ablation compares it with a Cartesian product of independent intervals
for $q$ and $\kappa_3$. The results in
Section~\ref{sec:results} measures the remaining over-enclosure.

\section{Certified continuous separation}
\label{sec:separation}

This section turns point-risk and curvature enclosures into a global upper
bound.  The final verifier uses direct mixture-ratio interval boxes rather
than the specialized common-mean formula, providing formula-level
independence.

\begin{lemma}[Semiconvex secant majorant]
\label{lem:secant}
Let $f$ be twice differentiable on $[a,b]$ and $f''\ge-K$ there, with $K\ge0$.
For $x=a+s(b-a)$, $s\in[0,1]$,
\[
 f(x)\le(1-s)f(a)+sf(b)+\frac{K}{2}(b-a)^2s(1-s).
 \label{eq:secant}
\]
Consequently, the cell maximum is bounded by the maximum of the two endpoints
and, when its vertex lies in $[0,1]$, the vertex value of the quadratic on the
right.
\end{lemma}

\begin{proof}
The function $f(x)+Kx^2/2$ is convex.  Apply the defining secant inequality
for a convex function and collect the quadratic terms.  No assumption about
the sign of $f'$ is required.
\end{proof}

\begin{proposition}[Risk quadrature enclosure]
\label{prop:quadrature}
Fix $t\in[-m,m]$ and $A>0$.  Composite Simpson quadrature on $[-A,A]$,
with an interval enclosure of the fourth derivative on every panel, encloses
the truncated risk integral.  The omitted tails satisfy
\[
 0\le\int_{|z|>A}q(t,z)^2\phiN(z)\,dz
 \le 8m^2\frac{\phiN(A)}{A}.
\]
\end{proposition}

\begin{proof}
The panel error is $-h^5H^{(4)}(\xi)/2880$ when $h$ denotes the full panel
width.  Replacing $H^{(4)}(\xi)$ by its interval range is inclusion
monotone.  Since $|\mu-t|\le2m$, the tail is at most
$4m^2\Pr(|Z|>A)$; Mills' inequality
$\Pr(|Z|>A)\le2\phiN(A)/A$ gives the result.
\end{proof}

The verifier obtains a lower enclosure for \eqref{eq:curvature} by partitioning
$[-A,A]$ in $z$, evaluating the common posterior ratio on the box
$t\in[a,b]$, and adding a two-sided curvature-tail enclosure.  If
$0\le V\le m^2$, then
\[
 |(V-1)^2+q\kappa_3|
 \le \max\{1,(m^2-1)^2\}+4m^4,
\]
because $|q|\le2m$ and
$|\kappa_3|\le2mV\le2m^3$.  This gives an explicit tail bound for
\eqref{eq:curvature}.

\begin{theorem}[Verifier implication]
\label{thm:verifier}
Suppose an outward-rounded computation returns bounds satisfying
\[
 B_L\le\BR(\pi)
 \quad\text{and}\quad
 U_U\ge\sup_{t\in[-m,m]}\risk(t,\delta_\pi).
\]
Then
\[
 \vstar\in[B_L,U_U],\qquad
 0\le\vstar-B_L\le U_U-B_L.
\]
If $U_U-B_L\le\eps$, the Bayes rule is certifiably $\eps$-minimax in value.
\end{theorem}

\begin{proof}
This is Corollary~\ref{cor:certificate} with interval endpoints weakened
outward.  It turns the established minimax bracket into an executable
certification rule.
\end{proof}

\section{Adaptive computational procedure}
\label{sec:algorithm}

The outer loop below is Kempthorne's algorithm specialized to symmetric
bounded-normal priors.  Symmetry reduces the search to nonnegative levels;
a positive level represents an equal pair at $\pm t$, while level zero
represents one atom.

\begin{figure}[ht]
\centering
\fbox{\begin{minipage}{0.91\linewidth}
\textbf{Adaptive support computation followed by certification.}
Given $m>0$, a requested gap $\eps>0$, and an initial symmetric prior:
\begin{enumerate}
\item reoptimize retained level masses;
\item optionally optimize interior support locations while keeping them ordered;
\item prune zero masses and merge numerically coincident levels;
\item compute a continuous global risk maximizer or upper bound on $[0,m]$;
\item if the numerical separation gap is not small, add the symmetric pair at
the maximizing level and return to step 1;
\item serialize only $m$, levels, masses, verifier settings, and claims;
\item run the independent outward-rounded verifier; and
\item return the prior, Bayes rule, $[B_L,U_U]$, and report only if the
verifier proves $U_U-B_L\le\eps$.  Otherwise refine and repeat.
\end{enumerate}
\end{minipage}}
\caption{The outer loop is inherited from Kempthorne; certification is a
separate post-processing stage.}
\label{alg:exchange}
\end{figure}

\paragraph{Fast mode.}
The discovery solver uses float64, deterministic Gauss--Hermite quadrature,
warm starts, quasi-Newton-style mass and location updates, aggressive pruning,
and the common-mean curvature enclosure.  Its bounds are diagnostics, not
certificates.  No parameter grid is fixed in advance; branch cells are
created adaptively.

\paragraph{Certificate mode.}
The final computation uses 192-bit Arb balls.  The verifier recomputes
posterior ratios, Simpson sums and remainder intervals, tail bounds, cell
curvatures, and continuous branch coverage.  It accepts no discovery-solver
branch tree.

\paragraph{Convergence statement.}
Kempthorne \cite[Theorem 3.1]{kempthorne1987} proves asymptotic convergence
under assumptions A1--A5 for the abstract exchange procedure.  The executable
guarantee here is conditional and post hoc: whenever the independent verifier
accepts, Theorem~\ref{thm:verifier} proves the reported gap.  A finite
iteration bound for finding a certifiable prior remains open.

\section{Independent certificate and verifier}
\label{sec:verifier}

The compact certificate contains $m$, nonnegative support levels, total level
masses, arithmetic precision, quadrature and branch settings, and claimed
risk intervals.  It contains no estimator samples, parameter grid, Monte
Carlo draws, or trusted branch partition.

The verifier performs the following checks.

\begin{enumerate}
\item It rejects nonfinite input, support outside $[0,m]$, unordered or
duplicate levels, nonpositive masses, and masses that do not sum exactly to
one as decimal inputs.
\item It reconstructs both atoms at every positive level and computes the
posterior directly from
$w_j\exp(\theta_jy-\theta_j^2/2)$.
\item It evaluates the risks at all support levels with independent composite
Simpson code, fourth-derivative intervals, and Gaussian tails, then averages
them to enclose the Bayes risk.
\item It starts from the complete interval $[0,m]$.  Each branch cell receives
a curvature lower bound and the exact semiconvex quadratic majorant.
Only a cell whose upper bound is within the work tolerance of the incumbent
may be pruned.
\item It checks that the supplied claim intervals contain the recomputation
and that the final gap is below the requested tolerance.
\end{enumerate}

The verifier imports neither \code{solver\_fast.py} nor
\code{solver\_certified.py}.  It does share the mathematical specification:
Gaussian density, posterior mean, risk, and the semiconvex theorem.  Complete
formula independence would be impossible without changing the quantity being
verified.  The implementation of those formulas and the arithmetic library
are independent.

\paragraph{Trusted computing base.}
The final proof trusts CPython, python-flint 0.8.0, the Arb/FLINT
implementation beneath it, and the verifier source.  The discovery optimizer
and NumPy are outside the final trusted base.  The resulting assurance level
is a computer-assisted proof with an explicit software trusted base.

\subsection{AI assistance}
AI tools assisted with literature search, theorem exploration, drafting, code
generation and adversarial testing in the project. During preparation of
this revision, OpenAI GPT-6.1 Terra assisted with manuscript revision and
GPT-6.1 Sol assisted with code and build work. AI-generated suggestions are not mathematical evidence. The verifier and its arithmetic
library form part of the trusted computing base; an independently executable
program is not an independent human review.

\section{Numerical results and ablations}
\label{sec:results}

All timings in this section are single-process measurements on the local
macOS 15.7.3 arm64 machine and are descriptive rather than hardware-normalized
complexity claims.

\subsection{The certified result for m=4}

The final prior has nonnegative levels
\[
 (0.5984699814893856,\;2.1264117002421385,\;4)
\]
and total level masses
\[
 (0.36699366016033996,\;0.3994285510645993,\;
   0.23357778877506074).
\]
Each positive level is split equally between its positive and negative atom.
The independent 192-bit verification used tail cutoff $A=8$, 2048 Simpson
panels per point risk, 512 curvature cells, and 567 branch nodes.  It took
49.68 seconds in the recorded clean run.

\begin{table}[ht]
\centering
\caption{Outward-rounded $m=4$ certificate.  Displayed endpoints are rounded
outward from the longer Arb output.}
\label{tab:m4}
\small
\begin{tabular}{ll}
\toprule
Quantity & Independently recomputed enclosure \\
\midrule
Bayes risk & $[0.8150887565050,\,0.8150887565548]$ \\
Worst-case risk & $[0.8150887565301,\,0.8150887615302]$ \\
Minimax risk & $[0.8150887565050,\,0.8150887615302]$ \\
Certified gap upper bound & $5.02517E-9$ \\
Branch nodes & 567 \\
Arithmetic & python-flint Arb balls, 192 bits \\
\bottomrule
\end{tabular}

\end{table}

The verifier's Bayes-risk width is $4.98\times10^{-11}$.  The branch envelope
uses a $5.0\times10^{-9}$ work width.  The certified minimax-gap upper bound is
$5.02517\times10^{-9}$, below $10^{-8}$.

\subsection{Fast deterministic sweep}

Table~\ref{tab:sweep} reports float64 discovery results.  They are not
certificates except where a separate certificate is stated.

\begin{table}[ht]
\centering
\caption{Adaptive-support float64 sweep; no ex ante parameter grid.}
\label{tab:sweep}
\small
\begin{tabular}{rrrrrr}
\toprule
$m$ & levels & atoms & Bayes risk & float gap & seconds \\
\midrule
0.25 & 1 & 2 & 0.0588158765 & 0.00e+00 & 0.003 \\
0.5 & 1 & 2 & 0.1989864336 & 0.00e+00 & 0.000 \\
1 & 1 & 2 & 0.4495995092 & 0.00e+00 & 0.001 \\
1.5 & 2 & 3 & 0.5561720123 & 3.78e-10 & 0.020 \\
2 & 2 & 3 & 0.6448008729 & 3.47e-09 & 0.030 \\
3 & 3 & 5 & 0.7506254421 & 3.62e-09 & 0.076 \\
4 & 3 & 6 & 0.8150887565 & 2.64e-09 & 0.187 \\
5 & 4 & 7 & 0.8573009664 & 3.71e-09 & 0.218 \\
8 & 6 & 11 & 0.9231345345 & 6.26e-09 & 5.149 \\
\bottomrule
\end{tabular}

\end{table}

\subsection{Dependency ablation}

Table~\ref{tab:ablation} holds the prior and endpoint-risk evaluator fixed.
The two generic global bounds exhaust the 5001-node budget.  Natural interval,
one-noise affine, Cartesian posterior-flow, and common-mean flow bounds
terminate.  These rows are float64 diagnostics; the Arb calculation supplies
the certified result.

\begin{table}[ht]
\centering
\caption{$m=4$, target separator width $10^{-8}$.}
\label{tab:ablation}
\small
\begin{tabular}{lrrrr}
\toprule
Separator & nodes & exhausted & width & seconds \\
\midrule
global Lipschitz & 5001 & yes & 2.34e-01 & 0.115 \\
global semiconvex & 5001 & yes & 5.09e-04 & 0.113 \\
natural interval & 291 & no & 8.58e-09 & 2.248 \\
one-noise affine & 427 & no & 9.95e-09 & 2.208 \\
Cartesian flow & 281 & no & 5.78e-09 & 3.170 \\
common-mean flow & 277 & no & 4.57e-09 & 0.076 \\
\bottomrule
\end{tabular}

\end{table}

The common-mean formula reduces the Cartesian-flow tree from 281 to 277 nodes.
Its much shorter runtime in this implementation comes primarily from
vectorized engineering.  Across cell
widths from approximately $3.87$ to $10^{-5}$, the ratio of common-mean to
Cartesian curvature width ranged from $0.864$ to $0.996$.  The measured
mathematical tightening is therefore modest.

\subsection{Numerical checks}

A seeded set of 100 valid symmetric priors and parameter intervals produced
no dense-sample violation of the float64 curvature enclosure. Such checks
provide falsification evidence; outward-rounded inclusion supplies the proof.

\section{Comparison with Montiel Olea and Zubova}
\label{sec:comparison}

Montiel Olea and Zubova \cite{montielzubova2026} use stochastic mirror ascent
on an equally spaced grid.  Their theory controls discretization and
stochastic optimization error, proves Wasserstein convergence of approximate
priors, and proves convergence of the resulting estimators.  The method is
expressed through broadly applicable first-order stochastic optimization.
The present solver exploits scalar Gaussian posterior identities to produce a
deterministic certificate for this bounded-normal instance.

\begin{table}[ht]
\centering
\caption{Complementary computational objectives.}
\label{tab:moz-comparison}
\small
\begin{tabular}{p{0.29\linewidth}p{0.31\linewidth}p{0.31\linewidth}}
\toprule
Feature & Montiel Olea--Zubova & Present project \\
\midrule
Parameter representation & Fixed equally spaced grid &
Adaptive continuous branch cells; finite active support\\
Optimization & Stochastic mirror ascent over grid weights &
Kempthorne exchange; float64 mass/location reoptimization\\
Integration & Monte Carlo supergradient and reported risk evaluation &
Deterministic Gauss--Hermite discovery; interval Simpson verification\\
Guarantee & High-probability approximation and convergence theory &
Post-hoc deterministic value certificate for each accepted prior\\
Continuous separation & Lipschitz grid correction &
Outward-rounded branch-and-bound on $[0,m]$\\
Generality & First-order method with broader optimization relevance &
Specialized scalar bounded-normal implementation\\
\bottomrule
\end{tabular}
\end{table}

For a numerical cross-check, we reconstructed their stated tuning formulas and
digitized the published risk figure using a checksum-locked pixel calibration.
The deterministic values for $m=1,2,3,4$ lie within the displayed figure
intervals to the declared pixel error.

At $m=4$ the present float64 discovery solver took 0.25 seconds in the clean
reproduction, and the independent certificate took 49.68 seconds.  These
times describe the current implementation on the recorded machine.  The two
methods target different computational guarantees, so the comparison is made
through their outputs and assurance levels rather than a speed ratio.

The deterministic certificate independently validates the bounded-normal
minimax value used in the local-projection/vector-autoregression application
of Montiel Olea and Zubova.  Its scope is the scalar minimax ingredient.  The
stochastic method retains the broader optimization and estimator-convergence
theory established in their paper.

\section{Scope and conclusion}
\label{sec:conclusion}

The final certificate concerns a one-dimensional normal location model with
unit variance, squared loss, bounded symmetric parameter space, and a supplied
finite symmetric prior.  Kempthorne's convergence theorem governs the abstract
exchange procedure; the executable guarantee here is post hoc: acceptance by
the independent verifier proves the reported minimax bracket.  The specialized
branch tree has no finite worst-case complexity bound in this release.  Its
trusted base is CPython, python-flint, Arb/FLINT and the verifier source.  The
$m=4$ result receives the final Arb certificate; the remaining sweep records
float64 discovery evidence.

The paper combines established least-favorable-prior optimization with a
continuous separator and independent outward-rounded verification.  The
common-mean enclosure gives a modest node reduction in the reported ablation;
the compact executable certificate is the main deliverable.  For $m=4$ it
proves a minimax-risk interval of width below $5.026\times10^{-9}$ and exposes
every numerical tolerance required to replay that statement.

\section*{Declarations}

\paragraph{Code and reproducibility.}
The implementation, final certificate, lockfile and reproduction commands are
available at
\href{https://github.com/Maar10Herr/Bounded-Normal-Minimax-Certificate-Public}
{github.com/Maar10Herr/Bounded-Normal-Minimax-Certificate-Public}.

\paragraph{Software and hardware.}
The recorded environment used CPython 3.13.5, NumPy 2.3.2, Pillow 11.3.0 and
python-flint 0.8.0 on macOS 15.7.3 arm64.  Reported timings are descriptive
machine measurements; the certified inequalities follow from the outward-
rounded verifier.

\paragraph{Acknowledgements.}
The project was motivated by the bounded-normal minimax work of Jos\'e Luis
Montiel Olea and Ekaterina Zubova; their work and the classical literature
are cited throughout.


\appendix
\section{Proof details}
\label{app:proofs}

\subsection{Differentiation of Bayes risk}

Fix $\alpha_0<1$.  For $\alpha\in[0,\alpha_0]$,
$f_{0,\alpha}\ge(1-\alpha_0)f_0$.  Directly expanding the differentiated
integrand in Proposition~\ref{prop:direction} reduces it to the sum of
$-f_2+f_1^2/f_0$ and
$(t-\delta_\pi)^2\phiN(\,\cdot-t)$.  The first absolute value is bounded by
$f_2+m^2f_0$, and the second by $4m^2\phiN(\,\cdot-t)$.  Both are integrable.
Dominated convergence therefore justifies the right derivative.  The same
argument applies to a direction given by any bounded-support probability
measure after integrating the atom formula in $t$.

\subsection{Fourth derivative used by Simpson's rule}

Let $u(z)=q(t,z)^2$.  With derivatives taken in $z$,
\begin{align*}
u_0&=q^2,\\
u_1&=2qV,\\
u_2&=2(V^2+q\kappa_3),\\
u_3&=2(3V\kappa_3+q\kappa_4),\\
u_4&=2(3\kappa_3^2+4V\kappa_4+q\kappa_5).
\end{align*}
Because $\phiN'=-z\phiN$, repeated product differentiation gives
\[
(u\phiN)^{(4)}
=\phiN\!\left[
u_4-4zu_3+6(z^2-1)u_2+4(3z-z^3)u_1
 +(z^4-6z^2+3)u_0\right].
\]
Every posterior quantity is evaluated on the same $(t,z)$ box.

\subsection{Curvature tail}

For posterior support in $[-m,m]$, $0\le V\le m^2$ by the
Bhatia--Davis inequality.  Also
\[
 |\kappa_3|=\left|\E(\Theta-\mu)^3\right|
 \le 2m\,\E(\Theta-\mu)^2=2mV.
\]
Thus the absolute curvature state is no greater than
\[
C_m=\max\{1,(m^2-1)^2\}+4m^4.
\]
Outside $[-A,A]$, the magnitude omitted from \eqref{eq:curvature} is at most
$2C_m\Pr(|Z|>A)\le4C_m\phiN(A)/A$.

\subsection{Branch pruning invariant}

At every time the active heap plus already pruned cells covers $[0,m]$.
An active cell has a proved upper bound.  A cell is pruned only if its bound
is no larger than the current incumbent plus the work tolerance.  Therefore
the global supremum is bounded by the maximum of the active-cell bounds and
the incumbent plus the work tolerance.  This maximum is exactly the value
reported by the verifier.  The use of floating-point heap priorities affects
only processing order; each retained mathematical bound is an Arb ball.

\section{Implementation and verification details}
\label{app:numerics}

\subsection{Error budget at m=4}

The independently recomputed Bayes-risk enclosure width is
$4.9777\times10^{-11}$.  Of this, the reported Simpson-remainder width is
$4.9696\times10^{-11}$, the Gaussian tail upper bound is
$8.084\times10^{-14}$, and arithmetic rounding in the Simpson sum is below
$3.1\times10^{-40}$.  At risk points encountered by separation, the largest
fourth-order remainder width is $6.5690\times10^{-11}$ and the largest tail
upper bound is $8.084\times10^{-14}$.  The branch envelope width is
$5.0\times10^{-9}$.

\subsection{Independent reconstruction}

The discovery solver computes posterior states with vectorized NumPy and a
specialized flow tube.  The final verifier instead:
\begin{enumerate}
\item constructs Arb balls from decimal strings;
\item evaluates each mixture exponential directly;
\item forms posterior central moments through order five;
\item implements a separate Simpson panel loop;
\item uses direct interval boxes for curvature rather than the common-mean
flow formula; and
\item rebuilds its own branch tree from $[0,m]$.
\end{enumerate}
The two paths share the mathematical definition of risk while using different
numerical representations and control flow.

\subsection{Failure modes examined}

Independent intervals for $f_0$, $f_1^2$, and $1/f_0$ repeatedly count the
same uncertainty and become particularly wide in the tails.  A global
Lipschitz bound scales only linearly with cell width and exhausted the
5001-node budget even at tolerance $10^{-4}$.  A global curvature bound also
exhausted that budget.  One-noise affine arithmetic helps on narrow cells but
was sometimes wider than the posterior-flow representation.  These empirical
results characterize the tested representations and tolerances.

\section{Detailed prior-art comparison}
\label{app:lineage}

Kempthorne's Section 3 begins with a finite prior, locally maximizes Bayes risk
over weights and locations, globally maximizes the Bayes-rule risk, stops when
maximum risk equals Bayes risk, and otherwise inserts the maximizer
\cite[pp.~175--177]{kempthorne1987}.  He also notes that zero weights or
coincident locations can reduce support.  His implementation brackets local
risk maxima by a parameter grid and refines them by golden-section search
\cite[p.~177]{kempthorne1987}.  The present verifier adds outward rounding
and complete continuous branch coverage to that optimization lineage.

Gourdin et al.\ terminate their Lipschitz global optimizer when the gap
between an upper bound and incumbent falls below a chosen tolerance
\cite[pp.~19--21]{gourdin1994}.  Their bounded-normal discussion fixes a
finite-dimensional prior form justified for a small range of $m$
\cite[pp.~25--26]{gourdin1994}.  The present solver combines adaptive support
generation with an interval-certified continuous risk oracle.

Dytso et al.\ \cite{dytso2022} give support-cardinality reductions under a
finite observation alphabet and use projected gradient ascent in numerical
examples.  Their stated finite-alphabet theorem addresses a different
observation space from the scalar Gaussian model.  Kent et al.\
\cite{kent2021} assume derivative and local linear-minimization oracles and
use stochastic subroutines.  Guggenberger and Huang
\cite{guggenberger2025} obtain iterationwise brackets in a
finite/discretized fictitious-play problem with grid-based best responses.
These results locate the current work within broader finite-mixture and
minimax-optimization methods.

Maruyama and Matsuda \cite{maruyama2025} provide the closest located
interval-assisted normal-mean minimaxity precedent.  They verify inequalities
for a Bayes estimator under a fixed half-Cauchy prior in a multivariate
problem.  The certified object here is instead the Bayes/maximum-risk gap for
a least-favorable prior in the bounded scalar problem.

This comparison is source-grounded and records the closest antecedents used in
the paper.  Future revisions should incorporate any closer continuous,
outward-rounded bounded-normal certificate brought to the author's attention.

\section{Reproducibility}
\label{app:reproducibility}

The repository linked in the declarations contains the verifier, certificate,
lockfile and setup instructions. From its root, the certificate can be replayed
with
\begin{verbatim}
uv run --frozen python src/verification/verify_certificate_arb.py \
  certificates/final/m4_epsilon_1e-8_arb.json
\end{verbatim}
The verification reconstructs the Bayes-risk lower bound and the global
risk upper bound for the specified finite symmetric prior. Its conclusion is
conditional on the trusted computational base described in
Section~\ref{sec:verifier}.


\bibliographystyle{plain}
\bibliography{bibliography}

\end{document}
